\subsection{Ràng buộc của môi trường di động}
\label{sec:mobile-constraints}

\subsubsection{Mobile App so với Web App và Desktop App}
\label{sec:mobile-vs-desktop}

Yêu cầu L.O.1.1 của học phần yêu cầu phân tích sự khác biệt kiến trúc giữa ứng dụng di động và ứng dụng máy tính. Bảng~\ref{tab:mobile-desktop} so sánh ba loại ứng dụng theo các tiêu chí ảnh hưởng trực tiếp đến thiết kế của OneSpace.

\begin{longtable}{>{\raggedright\arraybackslash}p{2.6cm} >{\raggedright\arraybackslash}p{4.4cm} >{\raggedright\arraybackslash}p{4.2cm} >{\raggedright\arraybackslash}p{4.2cm}}
\caption{So sánh Native Mobile App, Web App và Desktop App}
\label{tab:mobile-desktop}\\
\toprule
\textbf{Tiêu chí} & \textbf{Native Mobile App} & \textbf{Web App} & \textbf{Desktop App} \\
\midrule
\endfirsthead
\toprule
\textbf{Tiêu chí} & \textbf{Native Mobile App} & \textbf{Web App} & \textbf{Desktop App} \\
\midrule
\endhead
\bottomrule
\endfoot
Hiệu năng & Tốt, truy cập API hệ điều hành trực tiếp & Phụ thuộc trình duyệt & Thường có nhiều tài nguyên hơn \\
RAM, CPU & Giới hạn theo thiết bị & Phụ thuộc trình duyệt và thiết bị & Thường ít bị giới hạn \\
Pin & Cần tối ưu nghiêm ngặt & Cũng cần tối ưu trên điện thoại & Ít bị ảnh hưởng khi cắm nguồn \\
Phát nhạc nền & MediaSession / Service & Có thể bị giới hạn bởi trình duyệt & Tương đối linh hoạt \\
Chạy game & Native rendering & Canvas / WebGL & Có thể dùng GPU mạnh \\
Offline & Room, DataStore, file cache & Service Worker, IndexedDB & Local filesystem / database \\
Vòng đời & Do hệ điều hành kiểm soát chặt & Phụ thuộc vòng đời tab/trình duyệt & Ít bị hạn chế hơn \\
Thông báo & Android Notification API & Web Notifications khi được hỗ trợ & OS notification \\
Phù hợp dự án & \textbf{Rất phù hợp} & Hạn chế ở tích hợp nền & Không phù hợp mục tiêu mobile \\
\end{longtable}

Kết luận: ứng dụng tiêu dùng nội dung giải trí đa phương tiện cần phát nhạc nền ổn định, xử lý vòng đời nghiêm ngặt và thông báo tức thì~--~đây là những khả năng chỉ native mobile app đáp ứng đủ cho phạm vi môn học.

\subsubsection{Phân hệ nào tiêu thụ nhiều tài nguyên}
\label{sec:resource-consumption}

Bảng~\ref{tab:resource-consumption} đánh giá tương đối mức tiêu thụ tài nguyên theo phân hệ (định tính, không phải benchmark), làm cơ sở cho chính sách tối ưu ở mục~\ref{sec:resource-policy}.

\begin{longtable}{>{\raggedright\arraybackslash}p{1.6cm} >{\centering\arraybackslash}p{1.3cm} >{\centering\arraybackslash}p{1.8cm} >{\centering\arraybackslash}p{1.8cm} >{\raggedright\arraybackslash}p{8.0cm}}
\caption{Mức tiêu thụ tài nguyên theo phân hệ (đánh giá định tính)}
\label{tab:resource-consumption}\\
\toprule
\textbf{Phân hệ} & \textbf{Pin} & \textbf{Dữ liệu mạng} & \textbf{CPU/RAM} & \textbf{Biện pháp} \\
\midrule
\endfirsthead
\toprule
\textbf{Phân hệ} & \textbf{Pin} & \textbf{Dữ liệu} & \textbf{CPU/RAM} & \textbf{Biện pháp} \\
\midrule
\endhead
\bottomrule
\endfoot
Video & Cao & Rất cao & Cao & Adaptive bitrate, giới hạn preloading, giải mã phần cứng \\
Music & Thấp--TB & Thấp--TB & Thấp & Audio cache, phát với màn hình tắt, tránh polling \\
News & Thấp & Thấp--TB & Thấp & Cache bài viết, lazy loading ảnh \\
Games & TB--Cao & Thấp nếu offline & Phụ thuộc game & Giảm FPS, dừng game loop khi pause \\
\end{longtable}

Video thường tiêu thụ dữ liệu di động nhiều nhất và là một trong những nguồn hao pin lớn nhất. Tuy nhiên, game có đồ họa phức tạp vẫn có thể tốn pin hơn video~--~đây là lý do feature F-GAME-15 (kiểu flappy) được chấm Effort 2 thay vì 1 tại bảng~\ref{tab:feature-pool}.

\subsubsection{Chính sách tối ưu tài nguyên}
\label{sec:resource-policy}

Nhóm xây dựng một \texttt{ResourcePolicyManager} điều phối các quy tắc xử lý theo tình huống:

\begin{longtable}{>{\raggedright\arraybackslash}p{5.2cm} >{\raggedright\arraybackslash}p{9.8cm}}
\caption{Chính sách xử lý tài nguyên theo tình huống}
\label{tab:resource-policy}\\
\toprule
\textbf{Tình huống} & \textbf{Quy tắc xử lý} \\
\midrule
\endfirsthead
\toprule
\textbf{Tình huống} & \textbf{Quy tắc xử lý} \\
\midrule
\endhead
\bottomrule
\endfoot
Video chuyển thành Floating & Chỉ giữ một video decoder đang hoạt động \\
Music chạy nền & Dùng MediaSessionService, không giữ UI chạy \\
Game chuyển tab & Pause game loop, lưu game state \\
News ngoài màn hình & Ngừng các tác vụ UI không cần thiết \\
RAM thấp & Giảm prefetch, giải phóng player và ảnh không dùng \\
Bật Data Saver & Giảm chất lượng streaming, hạn chế tải trước \\
Battery Saver & Giảm xử lý nền không thiết yếu \\
\end{longtable}

Đặc biệt, khi game không còn hiển thị, không giữ vòng lặp render 60~FPS chỉ để người dùng có thể quay lại nhanh~--~điều này gây lãng phí trên thiết bị RAM 3~GB của persona Quốc Huy (mục~\ref{sec:persona-quochuy}).

\subsubsection{Offline và mạng không ổn định}
\label{sec:offline-design}

Bảng~\ref{tab:offline-behavior} định nghĩa hành vi từng phân hệ khi mất mạng và khi có mạng trở lại:

\begin{longtable}{>{\raggedright\arraybackslash}p{1.6cm} >{\raggedright\arraybackslash}p{6.8cm} >{\raggedright\arraybackslash}p{6.8cm}}
\caption{Hành vi offline theo phân hệ}
\label{tab:offline-behavior}\\
\toprule
\textbf{Phân hệ} & \textbf{Khi mất mạng} & \textbf{Khi có mạng trở lại} \\
\midrule
\endfirsthead
\toprule
\textbf{Phân hệ} & \textbf{Khi mất mạng} & \textbf{Khi có mạng trở lại} \\
\midrule
\endhead
\bottomrule
\endfoot
Video & Phát phần còn buffer/cache; nếu không đủ thì chờ mạng & Retry hoặc resume từ vị trí trước đó \\
Music & Phát bài đã tải xuống; streaming chỉ phát được phần buffer & Tiếp tục streaming \\
News & Đọc bài đã cache, bookmark vẫn hoạt động & Cập nhật nội dung từ server \\
Games & Game offline vẫn chơi được (F-GAME-12/13) & Đồng bộ điểm số nếu có tài khoản \\
\end{longtable}

Lưu ý quan trọng: \textbf{cache không đồng nghĩa với download}. Streaming cache có thể bị xóa và không bảo đảm phát ngoại tuyến toàn bộ nội dung. Nếu muốn chức năng Downloaded Music (F-MUS-12), cần thiết kế lưu trữ và quản lý file riêng~--~tính năng này đã được xếp vào nhóm Thankless (hoãn v1) tại bảng~\ref{tab:feature-pool}.

\subsubsection{Cấu hình tối thiểu hỗ trợ}
\label{sec:min-config}

\begin{longtable}{>{\raggedright\arraybackslash}p{2.6cm} >{\raggedright\arraybackslash}p{12.4cm}}
\caption{Cấu hình tối thiểu đề xuất}
\label{tab:min-config}\\
\toprule
\textbf{Thành phần} & \textbf{Quyết định} \\
\midrule
\endhead
\bottomrule
\endfoot
Android & 8.0+ (API 26+) \\
RAM & 3~GB mục tiêu (khớp persona Quốc Huy, thiết bị phổ thông) \\
Storage & Khoảng 200~MB trống cho cài đặt và cache MVP \\
Internet & Cần cho streaming và news mới \\
Offline & News cached, downloaded music, minigame offline \\
iOS & \textbf{Không hỗ trợ trong MVP}~--~phân tích iOS có ở mục~\ref{sec:platform-selection} nhưng nhóm không code bản iOS \\
\end{longtable}

API 26 và RAM 3~GB là mục tiêu hỗ trợ đề xuất, không phải mức tối thiểu bắt buộc của Jetpack Compose. Cần kiểm thử trên emulator hoặc thiết bị tương ứng (danh sách test case ở mục~\ref{sec:lifecycle-state}).
